\begin{tikzpicture}[
  x=1pt,y=-1pt,
  line width=.46pt,
  annotate/.style={fill=white,align=center,inner sep=1.7pt},
  participant/.style={draw,fill=white,align=center,minimum height=27pt,inner sep=2pt,font=\bfseries},
  activation/.style={draw,fill=white},
  lifeline/.style={dash pattern=on 3pt off 2pt},
  call/.style={-{Stealth[length=4.7pt,width=4.4pt]}},
  reply/.style={dash pattern=on 3pt off 2pt,-{Straight Barb[length=4.4pt,width=4.4pt]}},
  every node/.style={font=\linespread{1}\fontsize{10bp}{11.4bp}\selectfont}
]
  \foreach \x in {27,91,157,236,307,378} {
    \draw[lifeline] (\x,27) -- (\x,304);
  }
  \node[participant,minimum width=49pt,anchor=north] at (27,0) {Web 客户端};
  \node[participant,minimum width=60pt,anchor=north] at (91,0) {\texttt{AiChat}\\\texttt{Controller}};
  \node[participant,minimum width=65pt,anchor=north] at (157,0) {\texttt{AiChatService}\\\texttt{Impl}};
  \node[participant,minimum width=64pt,anchor=north] at (236,0) {\texttt{DeepSeek}\\\texttt{ApiClient}};
  \node[participant,minimum width=63pt,anchor=north] at (307,0) {\texttt{LocalAi}\\\texttt{Responder}};
  \node[participant,minimum width=71pt,anchor=north] at (378,0) {\texttt{AiChatHistory}\\\texttt{Service}};

  \draw (104,109) rectangle (417,161);
  \draw (104,109) -- (135,109) -- (135,119) -- (130,124) -- (104,124);
  \node[font=\bfseries] at (119,116.5) {opt};
  \node[annotate,anchor=west] at (166,118) {[配置了非空白 API 密钥]};
  \draw (104,170) rectangle (417,222);
  \draw (104,170) -- (135,170) -- (135,180) -- (130,185) -- (104,185);
  \node[font=\bfseries] at (119,177.5) {opt};
  \node[annotate,anchor=west] at (166,179) {[\texttt{message == null}]};

  \draw[activation] (87.5,45) rectangle (94.5,294);
  \draw[activation] (153.5,64) rectangle (160.5,276);
  \draw[activation] (374.5,85) rectangle (381.5,100);
  \draw[activation] (232.5,138) rectangle (239.5,154);
  \draw[activation] (303.5,199) rectangle (310.5,214);
  \draw[activation] (374.5,242) rectangle (381.5,259);

  \draw[call] (27,45) -- node[annotate,above] {\texttt{POST /api/ai/chat}} (87.5,45);
  \draw[call] (94.5,64) -- node[annotate,above] {\texttt{chat(request)}} (153.5,64);
  \draw[call] (160.5,85) -- node[annotate,above] {\texttt{requireWritableSession(sessionId)}} (374.5,85);
  \draw[reply] (374.5,100) -- node[annotate,above] {会话写权限校验通过} (160.5,100);

  \draw[call] (160.5,138) -- node[annotate,above] {同步请求模型} (232.5,138);
  \draw[reply] (232.5,154) -- node[annotate,above] {响应或异常} (160.5,154);
  \node[annotate,anchor=west] at (248,140) {取第一候选的非空白正文};
  \node[annotate,anchor=west] at (248,153) {运行时异常时回复仍为空};

  \draw[call] (160.5,199) -- node[annotate,above] {\texttt{reply(request)}} (303.5,199);
  \draw[reply] (303.5,214) -- node[annotate,above] {本地规则回复} (160.5,214);

  \draw[call] (160.5,242) -- node[annotate,above] {\texttt{save(request, response)}} (374.5,242);
  \node[annotate,anchor=east] (saveNote) at (368,249.5)
    {\texttt{@Transactional} 内重校会话并写入记录};
  \draw[densely dotted] (saveNote.east) -- (374.5,249.5);
  \draw[reply] (374.5,259) -- node[annotate,below] {保存正常完成} (160.5,259);

  \draw[reply] (153.5,276) -- node[annotate,above] {回复对象} (94.5,276);
  \draw[reply] (87.5,294) -- node[annotate,above] {成功响应} (27,294);
\end{tikzpicture}
